\documentclass[10pt,a4paper]{amsart}
\usepackage[margin=19mm]{geometry}
\usepackage{amsmath,amssymb,mathtools}
\usepackage[scr=rsfs,cal=euler]{mathalfa}
\usepackage{xfrac}
\usepackage[unicode,hidelinks,bookmarks=false]{hyperref}
\newcommand{\ie}{i.e.\ }
\newcommand{\cf}{cf.\ }
\newcommand{\resp}{respectively\ }
\newcommand{\Subsection}{\subsection}
\providecommand{\nobreakdash}{\nobreak\hbox{-}}
\setlength{\emergencystretch}{2em}
\multlinegap0pt
\usepackage{amssymb}
\usepackage{tikz}
\usepackage{csquotes}
\usepackage{mathtools}
\usepackage{bbm}
\newtheorem{theorem}{Theorem}
\newcommand{\C}{{\mathbb C}}
\DeclareMathOperator{\Hilb}{Hilb}
\newcommand{\N}{{\mathbb N}}
\DeclareMathOperator{\Proj}{\mathrm{Proj}}
\DeclareMathOperator{\Spec}{Spec}
\def\Sym{{\mathrm{Sym}}}
\title{A Lie-theoretic generalization of some Hilbert schemes}
\author{Oscar Kivinen}
\date{Source: arXiv:2512.08532v2 (CC BY 4.0)}
\begin{document}
\maketitle

\noindent\emph{Source and licence.} A Lie-theoretic generalization of some Hilbert schemes. Version arXiv:2512.08532v2 (CC BY 4.0), distributed by arXiv under the Creative Commons Attribution 4.0 International licence, http://creativecommons.org/licenses/by/4.0/, as recorded in the arXiv licence metadata for this exact version and shown on the abstract page https://arxiv.org/abs/2512.08532v2. Full licence notice: \texttt{licenses/arxiv-2512.08532-CC-BY-4.0.txt}.\par\smallskip

\noindent\textit{Editorial scope.} Counted unit: the article's theorem thm:haimanhilb (source0.tex lines 256--259). The article's complete original proof of it is delivered with it (source0.tex lines 260--263, one \texttt{proof} environment of 4 lines). No mathematics is written, repaired or re-derived: the statement and the proof are the article's own lines, in the article's own notation, and no retained line is substituted or re-flowed.  No further source-local context is needed: every cross-reference of every retained line resolves inside the two delivered spans.  Nothing was omitted from any retained span, so the package carries no omission record.  No retained line contains a citation, so the wrapper carries no bibliography and no bibliography entry.  No retained line contains a figure environment, an image-inclusion command or drawing code, so no image is delivered and the package declares no asset dependency.  Editorial formatting only, in addition: an AMS \texttt{amsart} preamble that restates the article's own declarations and macros used by the retained lines, namely the environments \texttt{theorem} on separate counters, together with the standard packages those lines need.  The retained environments print in this document as Theorem 1 in order of appearance, whereas the article prints the same items with the numbering of its own full document.  The complete native source of the article as distributed in the arXiv e-print for this exact version is bundled unchanged as \texttt{sources/190660/source0.tex}.
\begin{theorem}
\label{thm:haimanhilb}
$$\Hilb^n(\C^2)=\Proj\bigoplus_{d\geq 0} (\Delta A)^d,$$ where $\Delta A\subseteq \C[x_1,\ldots,x_n, y_1,\ldots,y_n]^{S_n}$ is the ideal consisting of products of the form $\Delta f$, where $\Delta=\prod_{i<j}(x_i-x_j)$ and $f\in A:=\C[x_1,\ldots, x_n,y_1,\ldots,y_n]^{sgn}$ is an alternating polynomial for the diagonal action.
\end{theorem}

\begin{proof} (Sketch.)
    On $\Spec \C[x_1,\ldots,x_n, y_1,\ldots, y_n]^{S_n}=\Sym^n \C^2$ there is an open locus $V=(\Sym^n\C^2)^{reg}$ of collections of $n$ points where all the points are distinct. Restricted to $V$, the Hilbert-Chow map $\Hilb^n(\C^2)\to \Sym^n \C^2$ is an isomorphism.
    For every partition $\mu\vdash n$, there are open charts $U_\mu \subseteq \Hilb^n(\C^2)$ consisting of those colength $n$ ideals $I\subseteq \C[x,y]$ for which $\C[x,y]/I$ is spanned by the monomials corresponding to the diagram of $\mu \subset \N\times \N$. On $U_\mu\cap V$, one can construct regular functions $\Delta_D/\Delta_\mu$, where $\Delta_D:=\det(x_i^{p_j}y_i^{q_j})_{i,j=1}^n$ are alternating polynomials attached to $n$-element subsets $D=\{(p_1,q_1),\ldots, (p_n,q_n)\}\subset \N\times \N$. One proves that these regular functions glue as we vary $\mu$ and also extend to regular functions on all of $\Hilb^n(\C^2)$. The universal property of blowing up then implies that there is a surjective map $\alpha: \Hilb^n(\C^2)\to \Proj\bigoplus_{d\geq 0} (\Delta A)^d$. Further, one checks that the pullbacks of the above $\Delta_D/\Delta_\mu$ give all the regular functions on $U_\mu$, showing that the map on structure sheaves is also surjective and $\alpha$ must be an isomorphism.
\end{proof}

\end{document}
